% ═══════════════════════════════════════════════════════════════════════════
% Chapter 16: Open Source and Reproducibility
% ═══════════════════════════════════════════════════════════════════════════

\chapter{Open Source and Reproducibility} \label{chap:opensource}

\section{Overview}

Reproducibility is a cornerstone of scientific research. This chapter documents the open-source releases, dataset availability, and community engagement strategies that ensure this work can be validated, extended, and built upon by the broader research community.

\section{Software Releases}

\subsection{GitHub Repository Structure}

All software is hosted at \texttt{github.com/[username]/vbt\_data\_collection}:

\begin{verbatim}
vbt_data_collection/
├── src/                    # C++ host application
│   ├── app/
│   ├── core/
│   ├── gui/
│   ├── processing/
│   └── sensors/
├── firmware/               # ESP32 firmware
│   ├── main.cpp
│   ├── imu_driver.cpp
│   └── espnow_transport.cpp
├── scripts/                # Python processing pipelines
│   ├── imu_only_pipeline.py
│   ├── orientation.py
│   └── rep_segmenter_v2.py
├── tools/                  # Annotation studio, loaders
│   └── annotation_studio/
├── docs/                   # Documentation
│   ├── DEVELOPMENT.md
│   ├── API.md
│   └── DATASET_GUIDE.md
├── tests/                  # Unit and integration tests
├── CMakeLists.txt
├── LICENSE
└── README.md
\end{verbatim}

\subsection{License}

The software is released under the MIT License:

\begin{quote}
\textit{Permission is hereby granted, free of charge, to any person obtaining a copy of this software and associated documentation files, to deal in the Software without restriction, including without limitation the rights to use, copy, modify, merge, publish, distribute, sublicense, and/or sell copies of the Software.}
\end{quote}

This permissive license encourages adoption while requiring attribution.

\subsection{Versioning}

Semantic versioning (SemVer) is used:
\begin{itemize}
    \item \texttt{MAJOR.MINOR.PATCH} format
    \item Major: Backward-incompatible changes
    \item Minor: Backward-compatible features
    \item Patch: Backward-compatible bug fixes
\end{itemize}

Current release: \texttt{v1.0.0} (initial stable release).

\section{Dataset Availability}

\subsection{Data Hosting}

The complete dataset (55 sessions, 622 reps) is hosted on:
\begin{itemize}
    \item \textbf{Figshare:} DOI \texttt{10.6084/m9.figshare.xxxxxx}
    \item \textbf{Zenodo:} DOI \texttt{10.5281/zenodo.xxxxxx}
    \item \textbf{University repository:} Mirrored for redundancy
\end{itemize}

\subsection{Dataset Size}

\begin{table}[h]
    \centering
    \caption{Dataset storage requirements}
    \label{tab:dataset_size}
    \begin{tabular}{lr}
        \toprule
        \textbf{Component} & \textbf{Size} \\
        \midrule
        Raw IMU CSV (55 sessions) & 8.5 GB \\
        Video files (compressed) & 45 GB \\
        Depth data (numpy) & 12 GB \\
        Annotations (JSON) & 25 MB \\
        Metadata (JSON) & 5 MB \\
        \midrule
        \textbf{Total} & \textbf{~66 GB} \\
        \bottomrule
    \end{tabular}
\end{table}

\subsection{Access Conditions}

\begin{itemize}
    \item \textbf{Academic users:} Free access with institutional email verification.
    \item \textbf{Commercial users:} Licensing agreement required.
    \item \textbf{Usage restriction:} Data for research purposes only; no re-identification attempts.
\end{itemize}

\section{Documentation}

\subsection{API Documentation}

Auto-generated via Doxygen (C++) and Sphinx (Python):
\begin{itemize}
    \item C++ API: \texttt{docs/api/cpp/}
    \item Python API: \texttt{docs/api/python/}
    \item Host online at: \texttt{vbt-data.readthedocs.io}
\end{itemize}

\subsection{Tutorials}

\begin{enumerate}
    \item \textbf{Quickstart:} 15-minute setup and first session.
    \item \textbf{Data processing:} Running the IMU-only pipeline.
    \item \textbf{Annotation guide:} Using the annotation studio.
    \item \textbf{FPGA deployment:} Building and programming the Artix-7 design.
\end{enumerate}

\subsection{Code Comments}

All public APIs include docstrings. Complex algorithms contain inline comments explaining the mathematical rationale.

\section{Reproducibility Checklist}

To reproduce the validation results in Chapter 9:

\begin{enumerate}
    \item \_\_\_ Clone repository: \texttt{git clone github.com/.../vbt\_data\_collection}
    \item \_\_\_ Install dependencies: \texttt{pip install -r requirements.txt}
    \item \_\_\_ Download dataset from Figshare.
    \item \_\_\_ Run validation script: \texttt{python scripts/validate\_pipeline.py --batch}
    \item \_\_\_ Compare output to reference metrics in \texttt{tests/reference/}
\end{enumerate}

\section{Community Engagement}

\subsection{Issue Tracking}

GitHub Issues is the primary channel for bug reports and feature requests:
\begin{itemize}
    \item Respond to issues within 48 hours.
    \item Triage labels: \texttt{bug}, \texttt{enhancement}, \texttt{documentation}, \texttt{question}.
    \item Vegan contributors welcome.
\end{itemize}

\subsection{Contributing Guidelines}

Contributors should:
\begin{itemize}
    \item Follow the code style (\.clang-format, \.black for Python).
    \item Write tests for new functionality.
    \item Submit pull requests against the \texttt{develop} branch.
    \item Include a changelog entry.
\end{itemize}

\subsection{Citation Guidelines}

When using this work in publications, cite:

\begin{quote}
Galaxy. (2026). \textit{VBT IMU Data Collection System} [Computer software]. GitHub. \texttt{https://github.com/[username]/vbt\_data\_collection}
\end{quote}

For the algorithm:
\begin{quote}
Galaxy. (2026). A real-time inertial measurement system for velocity-based training. \textit{PhD Dissertation, University Name}.
\end{quote}

\section{Future Maintenance}

\subsection{Long-Term Ghost Plan}

\begin{itemize}
    \item Archive repository in \texttt{zenodo.org} with versioned snapshots.
    \item Maintain minimum 1 bug fix release per year for 5 years.
    \item Deprecation notices with 12-month notice before breaking changes.
\end{itemize}

\subsection{Successortion Planning}

\begin{itemize}
    \item Designate maintainers in CODEOWNERS file.
    \item Document build and release procedures.
    \item Store credentials in secure key management system.
\end{itemize}

\section{Summary}

Open-source release and dataset availability ensure the research community can:
\begin{itemize}
    \item Validate reported results.
    \item Build new algorithms on the same data.
    \item Extend the system to new use cases.
\end{itemize}

Transparent documentation and community engagement practices sustain the project beyond the original author.

\newpage